\documentclass[11pt,a4paper]{article}
\usepackage[utf8]{inputenc}
\usepackage[T1]{fontenc}
\usepackage{geometry}
\geometry{margin=1in}
\usepackage{hyperref}
\usepackage{amsmath}
\usepackage{graphicx}
\usepackage{booktabs}
\usepackage{enumitem}
\usepackage{xcolor}
\usepackage{microtype}

\title{Vœrynth OS 5.0.1:\\A Local-First Neural Interface for Estate and Superyacht Automation}
\author{Kiran Karthikeyan Achari \and Danny Sneham}
\date{\today}

% Helper macro for screenshot placeholders.
% When you have real screenshots, place PNG files in the repository root
% (e.g., voerynth-dashboard-overview.png) and replace the commented
% \includegraphics line inside this macro with the real image include.
\newcommand{\figplaceholder}[2]{%
  \begin{figure}[t]
    \centering
    % TODO: when screenshots are available, replace the framed box below with:
    % \includegraphics[width=\linewidth]{../#2.png}
    \fbox{\parbox[c][5cm][c]{0.9\linewidth}{\centering Screenshot placeholder: #1\\[4mm]\texttt{#2.png} in repo root}}
    \caption{#1}
    \label{fig:#2}
  \end{figure}%
}

\begin{document}

\maketitle
\tableofcontents
\newpage

\begin{abstract}
Vœrynth OS 5.0.1 is the Layer~5 ``Neural Interface'' of Vœrynth Système, a local-first automation stack designed for high-end estates and superyachts. It is implemented as a React 18 single-page application, delivered as a web app, Capacitor-based Android app, and Electron desktop shell, and connects directly to a local Home Assistant instance that acts as the state and integration layer. Rather than acting as the ``brain'' of the system, Vœrynth OS presents a high-level command deck optimized for owners and staff: dashboards, room and zone decks, device and entity management views, and voice-ready entry points.

This paper provides a detailed description of Vœrynth OS 5.0.1: the domain requirements that motivated its design; the surrounding five-layer Vœrynth Système architecture; the internal UI and data architecture of the application; the development and deployment workflow; and a comparative analysis of how the interface differs from and improves on conventional Home Assistant dashboards for the target domain. We conclude with security and safety considerations and outline future work.
\end{abstract}

\section{Introduction}
Luxury automation systems are often marketed on the basis of visual polish and feature lists, yet they frequently fail at the fundamentals: low latency, robustness under poor connectivity, privacy, and operational clarity for non-technical staff. These shortcomings are amplified in marine environments, where connectivity is intermittent and safety constraints are strict. Vœrynth Système is built around a simple promise: the property must remain elegant and functional even when the internet is irrelevant. Vœrynth OS, the subject of this paper, is the visible surface of that promise.

Concretely, Vœrynth OS is a React-based command deck (web, Android, desktop) that connects to a local automation backbone in which Home Assistant~\cite{home-assistant} serves as the state and integration engine (Layer~3). Above that, Vœrynth Core (Layer~4) implements local reasoning, constraint logic, and auditability. Vœrynth OS (Layer~5) exposes real-time state, control surfaces, and explanatory views (``what happened and why'') to owners and staff.

The guiding principle is an ``iceberg'' model: 90\% of the engineering is invisible (Layers 1--4), and the remaining 10\% is the interface. The interface should not invite micromanagement; it should mainly serve for confirmation, exception handling, and deliberate control.

\figplaceholder{High-level Vœrynth OS Layer~5 dashboard overview}{voerynth-dashboard-overview}

\section{Domain Requirements and Design Objectives}

Vœrynth OS targets a specific class of deployments:
\begin{itemize}[leftmargin=*]
  \item Private estates with mixed wired/wireless subsystems (KNX lighting, IP HVAC, security zones, AV).
  \item Superyachts and vessels where satellite and cellular connectivity are intermittent and expensive.
  \item Hybrid environments where owners expect consistent behavior and presentation across shore properties and vessels.
\end{itemize}

From these deployments, several requirements follow:
\begin{itemize}[leftmargin=*]
  \item \textbf{Local-first operation.} All primary control and monitoring must work with only the local network available. Cloud services, if any, are treated as optional enhancements.
  \item \textbf{Low-latency, predictable control.} Interactions such as changing lighting scenes or HVAC setpoints must feel LAN-fast and should not silently depend on cloud round-trips.
  \item \textbf{Separation of roles.} Owners, captains, chief engineers, and house managers have different tasks; the interface must support role-tuned views.
  \item \textbf{Safety boundaries.} Automation must not be able to override certified marine control loops (autopilot, propulsion, class-approved alarms). Comfort and orchestration are in-scope; safety interlocks remain elsewhere.
  \item \textbf{Operational clarity for staff.} Staff need fast answers to ``what is happening now?'' and ``what changed?'' rather than a wall of toggles.
  \item \textbf{Deployability.} The system must be buildable and deployable by standard tooling (Node.js, npm, Android Studio, Electron) without bespoke infrastructure.
\end{itemize}

These requirements informed the decision to treat Vœrynth OS as Layer~5 of a deeper system and to design its internal architecture around exception-first workflows.

\section{Vœrynth Système Architecture}

Vœrynth Système is organized into five layers:
\begin{itemize}[leftmargin=*]
  \item \textbf{Layer 1 --- Hard Metal:} Physical connectivity (KNX/DALI, marine CAN/NMEA, discrete I/O, proprietary gateways).
  \item \textbf{Layer 2 --- Translation:} Signal~K~\cite{signal-k} and middleware normalize vendor-specific protocols into machine-readable event and state streams (e.g., MQTT/JSON).
  \item \textbf{Layer 3 --- State Engine:} Home Assistant Core~\cite{home-assistant} provides an integration and state registry with WebSocket and REST APIs.
  \item \textbf{Layer 4 --- Vœrynth Core:} Local reasoning, ``stupidity filters'', constraint logic, policy evaluation, and audit trails. This layer is responsible for deciding \emph{what} should happen and for logging \emph{why}.
  \item \textbf{Layer 5 --- Vœrynth OS:} The Neural Interface App described here, responsible for presenting state, accepting intent, and exposing traces.
\end{itemize}

Layers~1--4 are responsible for correctness, safety, and state continuity; Layer~5 is responsible for clarity, legibility, and control ergonomics.

The communication flow between Vœrynth OS and the lower layers is intentionally simple: Vœrynth OS talks directly to Home Assistant over its WebSocket and REST APIs. Vœrynth Core and other back-end services integrate with Home Assistant so that, from the perspective of the UI, complex reasoning is reflected as changes in entity state, attributes, and derived entities.

\figplaceholder{Conceptual communication flow between Vœrynth OS, Home Assistant, Vœrynth Core, and hardware layers}{voerynth-architecture-flow}

\section{Functional Overview of Vœrynth OS}

At a high level, Vœrynth OS provides the following capabilities:
\begin{itemize}[leftmargin=*]
  \item A \textbf{dashboard view} with a header, quick mode selector, and a curated set of cards (neural interface summary, weather, transit, global radar, flight tracker, energy flow, residents, ecosystem status).
  \item \textbf{Room/zone and role decks} that group controls and information for specific spaces and user roles.
  \item \textbf{Devices and entities management} using Home Assistant-style join logic for areas, devices, and entities, exposed in a single combined view.
  \item \textbf{Settings and system views} mirroring key Home Assistant configuration surfaces (areas, automations, integrations, people, system health, add-ons) but styled and organized for the Vœrynth aesthetic.
  \item \textbf{Voice-ready surfaces}, including hooks to trigger Assist flows and optional wake-word support on Android.
  \item \textbf{Multi-platform delivery}: a Vite-based web app, a Capacitor Android app, and an Electron desktop shell for fixed installations.
\end{itemize}

In the remainder of this section we describe the major user-facing surfaces in more detail.

\subsection{Dashboard and Quick Modes}

The primary landing view is implemented by the \texttt{DashboardView} component. It composes a header, a quick mode selector, and a grid of informational and control cards:
\begin{itemize}[leftmargin=*]
  \item \textbf{Dashboard header} shows a greeting, the current time, and high-level context for the estate or vessel.
  \item \textbf{Quick mode selector} exposes global modes (for example, ``At Sea'', ``In Port'', ``Guests Aboard'', ``Night'') that can drive coordinated changes in lighting, HVAC, privacy, and other subsystems.
  \item \textbf{Neural interface card} surfaces key signals from Vœrynth Core and summarises system confidence about its current behavior.
  \item \textbf{Weather and transit cards} display local conditions and relevant movement information.
  \item \textbf{Energy flow and residents cards} visualize energy usage and occupancy at a glance.
  \item \textbf{Ecosystem card} summarizes the health of underlying integrations (e.g., whether specific Home Assistant integrations are connected and online).
\end{itemize}

The dashboard supports an ``edit mode'', controlled via props, that allows the layout and individual cards to be reconfigured without changing the underlying back-end.

\figplaceholder{Dashboard with header, quick modes, and curated status cards}{voerynth-dashboard-cards}

\subsection{Devices and Entities View}

The Devices \\ Entities view is designed as an operator-friendly counterpart to Home Assistant's own device and entity registries. It is implemented in \texttt{DevicesEntitiesView} and backed by the global Home Assistant store (\texttt{useHAStore}). The view:
\begin{itemize}[leftmargin=*]
  \item Uses Home Assistant's WebSocket subscription to maintain an up-to-date map of all entity states.
  \item Fetches area, device, and entity registries via the \texttt{config/area\_registry/list}, \texttt{config/device\_registry/list}, and \texttt{config/entity\_registry/list} WebSocket commands.
  \item Performs the same join logic as the official Home Assistant frontend to compute derived fields such as "area name" for each entity (direct area assignment or via associated device).
  \item Exposes three tabs (Devices, Entities, Areas) in a single screen with a shared search bar and consistent card layout.
  \item Provides inline edit modals for entities, devices, and areas, allowing staff to adjust names, assignments, and other metadata without leaving the context.
  \item Offers a debug panel that displays subscription status, counts of areas/devices/entities, and loading state for troubleshooting.
\end{itemize}

Compared to a standard Home Assistant dashboard, which typically separates these concepts across multiple configuration screens, the combined Devices \\ Entities view reduces context switching and presents a more legible operational picture for staff.

\figplaceholder{Devices \\ Entities view with Devices, Entities, and Areas tabs and debug panel}{voerynth-devices-entities}

\subsection{Settings and System Surfaces}

The \texttt{src/views/settings} directory contains several views that mirror important Home Assistant configuration areas while preserving the Vœrynth visual language. These include:
\begin{itemize}[leftmargin=*]
  \item \textbf{Areas, Devices \\ Services, and Entity detail views} for drilling into the configuration of specific spaces or devices.
  \item \textbf{Automations and Integrations views} that list and categorize higher-level logic and external systems.
  \item \textbf{People and System views} for human-centric configuration and system diagnostics.
  \item \textbf{Add-ons view} to display and manage extensions in the underlying Home Assistant instance.
\end{itemize}

These views align with Home Assistant's mental model while remaining clearly within the Layer~5 command deck rather than the Layer~3 administration UI.

\figplaceholder{Settings home view aggregating areas, automations, integrations, people, and system status}{voerynth-settings-home}

\subsection{Configuration and Onboarding}

When Vœrynth OS is first launched, the user is prompted to configure the Home Assistant URL (for example, \texttt{http://192.168.1.100:8123}) and a long-lived access token. These values are stored locally using browser or device-secure storage. The configuration surface also exposes:
\begin{itemize}[leftmargin=*]
  \item Accent color selection, backed by an \texttt{AccentColorProvider} context, which controls Tailwind CSS-based color choices across the UI.
  \item Connection status indicators (connected, reconnecting, error) derived from the Home Assistant context and the underlying WebSocket connection.
  \item Optional wake-word and voice-assist toggles (particularly on Android builds), which can activate hands-free entry points to Assist.
\end{itemize}

\figplaceholder{First-run configuration view showing Home Assistant URL, token entry, and connection status}{voerynth-onboarding-config}

For documentation and support, visit \href{https://voerynth.de}{voerynth.de}.

\section{Internal Architecture and Implementation}

This section describes how Vœrynth OS is implemented, focusing on the front-end stack, Home Assistant communication layer, state management, and multi-platform delivery.

\subsection{Frontend Stack}

The application uses the following stack:
\begin{itemize}[leftmargin=*]
  \item \textbf{React 18}~\cite{react} for component-based UI composition.
  \item \textbf{Vite 7}~\cite{vite} as the build tool and development server, providing fast hot module reloading and optimized production bundles.
  \item \textbf{Tailwind CSS 4}~\cite{tailwind} for utility-first styling, with a small set of custom design tokens for the Vœrynth aesthetic.
  \item \textbf{Zustand} for lightweight state management, particularly for the global Home Assistant store.
  \item \textbf{Lucide React} icons and custom brand icons for visual affordances.
\end{itemize}

The project structure mirrors this organization, with dedicated directories for components, views, contexts, stores, services, hooks, and utilities, as documented in the repository README.

\subsection{Home Assistant Communication Layer}

Communication with Home Assistant is handled by two cooperating modules:
\begin{itemize}[leftmargin=*]
  \item \textbf{HAConnection} (\texttt{HAConnection.js}) manages a raw WebSocket connection to the \texttt{/api/websocket} endpoint of the configured Home Assistant instance. It is responsible for connection lifecycle, reconnection, and low-level message handling.
  \item \textbf{haClient} (\texttt{haClient.js}) wraps the \texttt{home-assistant-js-websocket} library and exposes higher-level functions such as \texttt{connect}, \texttt{subscribeEntities}, \texttt{callWS}, and \texttt{callREST}, as well as error normalization for common failure modes (cannot connect, invalid auth, connection lost, invalid HTTPS--HTTP combinations).
\end{itemize}

The \texttt{connect} function constructs an auth object from the configured URL and token, establishes a WebSocket connection (using the appropriate \texttt{ws://} or \texttt{wss://} prefix), and returns a connection object. The \texttt{callREST} helper performs authenticated REST calls to endpoints such as \texttt{/api/states} or \texttt{/api/history/period}, attaching the bearer token and raising exceptions on HTTP error status codes.

\subsection{Global Home Assistant Store}

The \texttt{useHAStore} Zustand store is responsible for mirroring Home Assistant state and registries on the client. It maintains:
\begin{itemize}[leftmargin=*]
  \item A map of entity states keyed by entity id.
  \item Arrays for areas, devices, entity registry entries, and integration configuration entries.
  \item Derived indices such as \emph{areasById}, \emph{devicesById}, and \emph{entity registry by entity id} for efficient lookup.
  \item Flags for loading state, errors, and subscription status.
\end{itemize}

The store exposes methods to initialize itself from a snapshot of states, subscribe to ongoing updates, refresh the registries, and reset to a known baseline. The Devices \\ Entities view uses these abstractions to implement HA-style join logic while keeping UI components relatively simple.

\subsection{Context Providers and Cross-Cutting Concerns}

Two primary React contexts wrap the application:
\begin{itemize}[leftmargin=*]
  \item \textbf{HomeAssistantProvider} encapsulates connection details, exposes connection status and raw Home Assistant state, and provides a method to retrieve the current WebSocket connection for modules such as \texttt{haClient}.
  \item \textbf{AccentColorProvider} stores the chosen accent color (for example, cyan, emerald) and exposes it to components so that Tailwind CSS class combinations can be constructed dynamically (such as \texttt{text-cyan-400} or \texttt{bg-emerald-600}).
\end{itemize}

The top-level \texttt{App} component nests these providers and then renders \texttt{AppContent}, which in turn hosts the navigation, dashboard, and other views.

\subsection{Multi-Platform Delivery}

Vœrynth OS is delivered on multiple platforms from a single codebase:
\begin{itemize}[leftmargin=*]
  \item \textbf{Web app:} The default target, built with Vite into a static \texttt{dist/} directory that can be served by any HTTP server.
  \item \textbf{Android app:} A Capacitor 7~\cite{capacitor} project (\texttt{android/}) wraps the web app in a native shell. The Capacitor configuration specifies the app id (\texttt{com.voerynth.os}), app name, and web directory. Android Studio is used to produce signed APKs for deployment.
  \item \textbf{Electron desktop:} An \texttt{electron/} directory contains a minimal Electron main process that either loads the Vite dev server (in development) or the built \texttt{index.html} (in production), with a frameless, menu-less window suitable for kiosk-style installs.
\end{itemize}

The repository also includes a Tauri-based harness (\texttt{src-tauri/}) that can be used to explore an alternative desktop packaging strategy, though Electron is the default path for cross-platform desktop builds in this release.

\section{Building and Deploying Vœrynth OS}

This section provides a reproducible path from a fresh checkout of the repository to a running system.

\subsection{Prerequisites}

\begin{itemize}[leftmargin=*]
  \item Node.js 18 or later and npm.
  \item A Home Assistant instance reachable on the same network as the device running Vœrynth OS.
  \item A long-lived access token for Home Assistant, created from the user profile page.
  \item For Android builds: Android Studio and an appropriate SDK.
  \item For Electron builds: a working Node.js environment on the target desktop platform.
\end{itemize}

\subsection{Web Development Workflow}

To run the app in development mode:
\begin{enumerate}[leftmargin=*]
  \item Clone the repository and install dependencies:
  \begin{verbatim}
git clone https://github.com/kiranvenom1209/voerynth-os.git
cd voerynth-os
npm install
  \end{verbatim}
  \item Optionally create a \texttt{.env} file to pre-configure the Home Assistant URL and token:
  \begin{verbatim}
VITE_HA_URL=http://192.168.1.100:8123
VITE_HA_TOKEN=your_long_lived_token_here
  \end{verbatim}
  \item Start the Vite dev server:
  \begin{verbatim}
npm run dev
  \end{verbatim}
  and open \texttt{https://os.voerynth.de} (or local dev URL) in a browser.
\end{enumerate}

\subsection{Production Builds}

For a production web build:
\begin{verbatim}
npm run build
  \end{verbatim}
This produces a \texttt{dist/} directory that can be served by a local HTTP server (for example, \texttt{nginx} running on the vessel or estate network).

For Android:
\begin{verbatim}
npm run build
npx cap sync android
npx cap open android
  \end{verbatim}
The project can then be built and signed from Android Studio to produce an APK or app bundle.

For Electron desktop builds:
\begin{verbatim}
npm run electron:build
  \end{verbatim}
which produces platform-specific artifacts suitable for installation on the target operating system.

\section{Comparison with Standard Home Assistant Dashboards}

Home Assistant provides a powerful and flexible dashboard system (Lovelace), and Vœrynth OS is deliberately not a replacement for it. Instead, it focuses on the Layer~5 use case of being a curated command deck for a specific property and set of roles. This section summarizes key differences.

\subsection{Interaction Model}

In a standard Home Assistant dashboard, each view is configured via YAML or a GUI editor. The default interaction model is card-centric: individual entities or groups are placed on a grid. This offers high flexibility but can lead to cluttered interfaces and a proliferation of dashboards tailored to individual preferences.

Vœrynth OS instead presents a small number of opinionated views:
\begin{itemize}[leftmargin=*]
  \item A main dashboard that summarizes the state of the property or vessel in a single screen.
  \item Device and entity management views with HA-style join logic but a unified layout and search.
  \item Settings surfaces that group related configuration areas under a coherent visual language.
\end{itemize}

This opinionated approach reduces configuration burden and provides a common operational picture for all staff.

\figplaceholder{Side-by-side comparison: standard Home Assistant Lovelace dashboard vs Vœrynth OS dashboard}{voerynth-vs-ha-dashboard}

\subsection{Operational Workflows}

For typical operational tasks, Vœrynth OS changes the workflow compared to a generic HA dashboard:
\begin{itemize}[leftmargin=*]
  \item \textbf{Changing global state:} Instead of toggling multiple scenes or modes across dashboards, an operator can select a quick mode (for example, ``Guests Aboard'') that orchestrates changes across subsystems.
  \item \textbf{Investigating issues:} The Devices \\ Entities view and debug panel provide an immediate overview of areas, devices, entities, and subscription health, without digging through multiple configuration pages.
  \item \textbf{Staff training:} A consistent, branded interface with fewer exposed controls makes it easier to train new crew or house staff compared to user-specific HA dashboards.
\end{itemize}

\subsection{Local-First Behavior}

Both Home Assistant and Vœrynth OS are capable of local operation, but Vœrynth OS is explicitly designed to avoid cloud dependencies at the interface layer. There is no default telemetry, no cloud authentication layer, and no reliance on third-party dashboards. For estates and vessels where connectivity is expensive or tightly controlled, this predictability is a significant advantage.

\subsection{Limitations}

Vœrynth OS does not attempt to replace the full configurability of Lovelace dashboards. Power users may still prefer to construct custom dashboards within Home Assistant for niche use cases. Vœrynth OS is deliberately focused on the shared command-deck scenario.

\section{Security, Privacy, and Safety}

Security and privacy considerations include:
\begin{itemize}[leftmargin=*]
  \item Vœrynth OS assumes that the Home Assistant instance is reachable only over trusted local networks or VPN; it does not encourage exposing Home Assistant directly to the public internet.
  \item Authentication is performed via long-lived access tokens, which are stored using platform-appropriate secure storage mechanisms (for example, browser storage with appropriate origin restrictions or mobile secure storage APIs).
  \item The application sends no telemetry by default. Any diagnostic mechanisms intended for support are opt-in and designed to exclude sensitive state.
\end{itemize}

From a safety perspective, Vœrynth Système explicitly excludes certified marine control loops (autopilot, propulsion, class-approved alarm chains) from its actuation surface. Vœrynth OS operates as a comfort, monitoring, and orchestration layer; it may request changes to lighting, HVAC, media, and similar subsystems, but hard safety interlocks reside in Layers~1--3.

\section{Future Work and Research Directions}

Future work on Vœrynth OS and the broader Vœrynth Système stack includes:
\begin{itemize}[leftmargin=*]
  \item Deeper integration with Vœrynth Core (Layer~4) to expose structured reasoning traces, including confidence levels and revert paths, directly in the UI.
  \item Multi-user and role-based access control, allowing different decks and permissions for owners, staff, guests, and technical crew.
  \item Richer degraded-mode behavior when subsets of the stack are offline, including clear health indicators and read-only fallbacks.
  \item Fleet-oriented configuration management for owners operating multiple vessels or properties.
  \item Formal user studies comparing effectiveness and error rates between standard HA dashboards and the Vœrynth OS command deck in operational scenarios.
\end{itemize}

In conclusion, Vœrynth OS 5.0.1 demonstrates that a luxury command-deck interface for estates and superyachts can be built on local-first principles without sacrificing usability. By treating the UI as the visible tip of a layered, safety-conscious architecture rather than as the system's brain, Vœrynth Système aims to deliver environments that feel intelligent without turning the owner into a full-time pilot.

\bibliographystyle{plain}
\bibliography{references}

\end{document}
